\section{System Design}
\label{sec:systemdesign}

This section presents the high-level architecture of the proposed
cabinet prototype and the rationale behind the principal component
choices.

\subsection{Architectural Overview}

The prototype consists of four cooperating subsystems organised
around an inexpensive Arduino Uno microcontroller. The sensor node
reads temperature and humidity from a DHT11, gas-trend data from an
MQ135 module, and shelf weight through an HX711 24-bit
ADC~\cite{hx711datasheet} connected to four 50\,kg load cells in a
Wheatstone-bridge configuration. The sensor cycle has a nominal
period of three seconds and emits one JSON record per cycle on the
microcontroller's USB serial port.

A camera node, implemented on an AI Thinker ESP32-CAM, joins the
local Wi-Fi network and exposes two HTTP endpoints: \texttt{/capture}
on port 80 returns a single JPEG suitable for QR decoding, and
\texttt{/stream} on port 81 serves an MJPEG live feed for the mobile
dashboard. An LED strip mounted around the cabinet's inner edges is
switched by a MOSFET driven from the Arduino's D9 pin and provides
diffused illumination for reliable camera operation.

A Python USB-to-HTTP bridge translates the serial stream into a
CORS-enabled local HTTP endpoint at
\texttt{http://localhost:8787/sensors}. The bridge also accepts
\texttt{POST} requests on \texttt{/tare} and \texttt{/recalibrate\_gas}
and converts them into single-character commands that the Arduino
sketch interprets as runtime calibration triggers. A cross-platform
Flutter application, built for Android and the web, polls the bridge
once per second, fetches frames from the ESP32-CAM directly, decodes
QR codes from cabinet captures, maintains an in-memory product list
with expiration metadata, and surfaces rule-based alerts.

Decoupling sensing from networking is a deliberate design decision.
The ATmega328P controller used by the Arduino Uno is constrained to
2\,kB of RAM and lacks native wireless capability; the
USB-to-HTTP bridge is a lightweight alternative to fitting the
sensor node with a Wi-Fi shield, and it preserves the educational
accessibility of the Arduino firmware.

\subsection{Subsystem Selection}

\subsubsection{Microcontroller}
The Arduino Uno was selected for its mature DHT and HX711 community
libraries, its single-file Wiring/C++ workflow, and its low cost. The
absence of native wireless capability is offset at the bridge layer
without altering the sensor firmware.

\subsubsection{Weight subsystem}
The HX711 amplifier is widely deployed in low-cost weighing
applications~\cite{hx711datasheet}. The four load cells are arranged
under a 47\,$\times$\,36\,cm wooden shelf and connected to the Uno via
\texttt{DT} and \texttt{SCK} on D4 and D5. The firmware exposes a
runtime tare command (\texttt{t}/\texttt{T}) and applies a slow
baseline-drift compensator to absorb the small negative excursions
introduced by load-cell creep.

\subsubsection{Environmental sensing}
The DHT11 sensor provides $\pm 2^\circ\mathrm{C}$ and $\pm 5\,\%$
relative-humidity accuracy~\cite{adafruitdht}, which is sufficient
for detecting the temperature transient associated with
refrigerated items being placed in a closed cabinet at room
temperature. Transient \texttt{NaN} reads are discarded and the
previous valid sample is forwarded, a robustness pattern that is
standard practice for this sensor.

\subsubsection{Gas observation}
The MQ135 is sampled through analog pin A0 with an in-firmware
exponential moving average. Low-cost MQ135 modules exhibit
significant baseline drift during the first 30--60\,s after power-on,
which the firmware absorbs without claiming absolute concentration
measurements. The bridge exposes a runtime endpoint
(\texttt{/recalibrate\_gas}) that re-establishes the baseline on
demand.

\subsubsection{Camera node}
The AI Thinker variant of the ESP32-CAM was selected for its local
availability in the Turkish hobbyist market and the maturity of its
firmware ecosystem~\cite{esp32cam}. Hosting the capture and stream
endpoints on different ports enables independent quality settings:
the capture endpoint operates at 320\,$\times$\,240 with low
compression, while the stream endpoint operates at
640\,$\times$\,480 with moderate compression, preventing
single-request stalls during interleaved access.

\subsubsection{Mobile application}
Flutter was selected to support a single codebase targeting both
Android phones and a browser-based web view. The web build is
published from the project's \texttt{docs/} directory and is
designed to operate in a self-contained demonstration mode when no
sensor hardware is attached, which is helpful in presentations
without the full physical prototype.

\subsection{Data Exchange Formats}

A representative serial record produced by the Arduino is:

\begin{lstlisting}
{"temperature":27.6,"humidity":62.0,
 "gasValue":43,"weight":247.45,
 "riskScore":10}
\end{lstlisting}

The bridge timestamps each record with the host wall-clock and
exposes a slightly extended representation on \texttt{GET /sensors}:

\begin{lstlisting}
{"temperature":27.7,"humidity":63.0,
 "gasValue":42,"weight":0,
 "updatedAt":1781692341,"online":true}
\end{lstlisting}

The Flutter application interprets \texttt{online:false} as an
out-of-band reset condition; the dashboard substitutes em-dashes for
the affected sensor readings and the alert subsystem raises an
\emph{ESP32 Sensor Board Offline} notice.
